% Fixed texts, English. Override single ones for a game in override-en.tex (\renewcommand).
\usepackage[english]{babel}
\newcommand{\tManual}{User manual}
\newcommand{\tUserManual}{USER MANUAL}
\newcommand{\tForV}{for V9990}
\newcommand{\tForTheV}{for the Yamaha V9990}
\newcommand{\tCoverLine}{The MSX game on the V9990}
\newcommand{\tOriginalBy}{An original game by}
\newcommand{\tPortBy}{V9990 port by}
\newcommand{\tHomage}{This is an unofficial port made for study, preservation, and homage.}
\newcommand{\tWhatYouNeed}{What you need}
\newcommand{\tChipCaption}{The Yamaha V9990 video processor, the chip this version of the game draws with.}
\newcommand{\tItem}{Item}
\newcommand{\tRequirement}{Requirement}
\newcommand{\tComputer}{Computer}
\newcommand{\tComputerText}{\textbf{Any MSX1} with two cartridge slots, one for the game and one for the V9990. Tested on a Sony HB-10P (MSX1) and on an MSX turbo R; other MSX generations should work but have not been tested.}
\newcommand{\tMemory}{Memory}
\newcommand{\tMemoryText}{\textbf{\ramNeeded{} of RAM} in the main memory (all four pages).}
\newcommand{\tVideo}{Video}
\newcommand{\tVideoText}{A \textbf{V9990 cartridge} (Yamaha V9990 video). It is mandatory: the whole picture is drawn by the V9990. The computer's own video chip is kept blank on purpose.}
\newcommand{\tSoundRow}{Sound}
\newcommand{\tSoundText}{The MSX PSG is always used. An \textbf{SCC} (in the game cartridge or in another slot) is optional.}
\newcommand{\tGameRow}{Game}
\newcommand{\tGameText}{The \textbf{\gameTitle{} V9990 cartridge}: a \romKind{} (the file \texttt{\romFile}, written to a flash cartridge or equivalent).}
\newcommand{\tMonitor}{Monitor}
\newcommand{\tMonitorText}{A monitor or TV connected to the \textbf{V9990's video output}, not to the MSX's own video output.}
\newcommand{\tControlsRow}{Controls}
\newcommand{\tControlsText}{Keyboard, cursor keys, or a joystick in port 1.}
\newcommand{\tNoTape}{\textbf{No tape, no disk, no password.} The whole game is in the cartridge.}
\newcommand{\tSettingUp}{Setting up}
\newcommand{\tSetOff}{\textbf{Switch the MSX off.} Never insert or remove a cartridge with the power on.}
\newcommand{\tSetV}{Insert the \textbf{V9990 cartridge} in one of the cartridge slots.}
\newcommand{\tSetGame}{Insert the \textbf{\gameTitle{} cartridge} in the other slot. The order does not matter.}
\newcommand{\tSetMonitor}{Connect your monitor or TV to the \textbf{video output of the V9990}.}
\newcommand{\tSetJoy}{Plug a joystick into port 1 if you want to use one.}
\newcommand{\tSetOn}{Switch the MSX on. For a moment the MSX's own screen shows \texttt{V9990 found !}; then the picture appears on the V9990's monitor. If it says \texttt{V9990 NOT found!!!} the V9990 is missing or not seated, and the game stops there.}
\newcommand{\tBlackScreen}{\textbf{If the screen stays black} check, in this order: the monitor is on the V9990's output; the V9990 and the game are fully inserted; the computer has \ramNeeded{} of RAM; the game ROM is the \romKind{} build. On an emulator such as openMSX, enable the V9990 extension and load the ROM with the right mapper.}
\newcommand{\tPowerOnTitle}{From power-on to the game}
\newcommand{\tPowerOnCaption}{Left: the cover. Right: the title screen.}
\newcommand{\tControlsTitle}{Controls}
\newcommand{\tCreditsTitle}{5. Credits and rights}
\newcommand{\tIsAGameBy}{is a game by}
\newcommand{\tRightsText}{The game, its graphics, music and name belong to their rights holders. This is an unofficial port made for study, preservation and homage; it is not affiliated with them.}
\newcommand{\tVersionWord}{V9990 version}
\newcommand{\tMsxTrademark}{MSX is a trademark of MSX Licensing Corporation. The MSX logo is shown only to say which computers the game runs on}
\newcommand{\tChipNote}{V9990 is the name of the Yamaha video chip; the chip photo is by Yaca2671, CC BY-SA 3.0.}
\newcommand{\tSourceNote}{Source code available on GitHub. \textbf{Please improve the game if you feel so!}}
\newcommand{\tAppendix}{Appendix. Cartridge label}
\newcommand{\tStickerText}{A label for the cartridge, ready to print. At 100\,\% it is about 16\,cm wide; scale it in your printing program to fit the label area of your cartridge (the artwork has a 3:2 proportion).}
\newcommand{\tAction}{Action}
\newcommand{\tKeyboard}{Keyboard}
\newcommand{\tJoystick}{Joystick}
\newcommand{\tNo}{No.}
\newcommand{\tElement}{Element}
\newcommand{\tTellsYou}{What it tells you}
\newcommand{\tTips}{Tips.}
